\chapter{操作系统接口}
\label{CH:UNIX}

操作系统的职责是在多个程序之间共享计算机资源，并提供比硬件本身更为丰富和有用的服务。操作系统对底层硬件进行管理和抽象，例如，文字处理程序无需关心所使用的磁盘硬件类型。操作系统在多个程序之间共享硬件，使它们能够同时运行（或看起来同时运行）。此外，操作系统为程序之间的交互提供受控的方式，使它们能够共享数据或协同工作。

操作系统通过接口向用户程序提供服务。
\index{interface design}
设计一个好的接口实际上相当困难。一方面，我们希望接口简单而精炼，因为这样更容易保证实现的正确性；另一方面，我们又可能倾向于为应用程序提供许多复杂的功能。解决这一矛盾的关键在于：设计依赖少数几种机制的接口，通过这些机制的组合来提供广泛的通用性。

本书以一个具体的操作系统为例来阐释操作系统的概念。该操作系统名为
xv6，它提供了 Ken Thompson 与 Dennis Ritchie 的 Unix 操作系统~\cite{unix} 所引入的基本接口，同时借鉴了 Unix 的内部设计。Unix 提供了一套精炼的接口，其各种机制组合效果良好，具有出人意料的通用性。这套接口取得了巨大成功，以至于现代操作系统——BSD、Linux、macOS、Solaris，乃至在一定程度上的 Microsoft Windows——都具有类 Unix 的接口。
理解 xv6 是理解上述系统乃至许多其他系统的良好起点。

如
图~\ref{fig:os} 所示，
xv6 采用传统的
\indextext{内核}形式——
一种为运行中的程序提供服务的特殊程序。
每个运行中的程序称为
\indextext{进程}，
其内存中包含指令、数据和栈。指令实现
程序的计算逻辑。数据是计算所操作的变量。栈则组织程序的过程调用。
一台计算机通常运行多个进程，但只有一个
内核。

当进程需要调用内核服务时，它会调用
\indextext{系统调用}——
操作系统接口中的调用之一。
系统调用进入内核；
内核执行服务后返回。
因此，进程在
\indextext{用户空间}
和
\indextext{内核空间}之间交替执行。

如后续章节所详述，内核利用 CPU\footnote{
本书通常用\indextext{CPU}（中央处理器的缩写）来指代执行计算的硬件单元。其他文献（如 RISC-V 规范）也使用处理器、核心和硬件线程（hart）等术语代替 CPU。
} 提供的硬件保护机制，确保在用户空间执行的每个进程只能访问自己的内存。
内核在执行时拥有实现这些保护所需的硬件特权级，而用户程序则不具备这些特权。
当用户程序调用系统调用时，硬件
会提升特权级，并开始执行内核中预先安排好的函数。

\begin{figure}[t]
\center
\includegraphics[scale=0.5]{fig/os.pdf}
\caption{内核与两个用户进程。}
\label{fig:os}
\end{figure}

内核提供的系统调用集合就是用户程序所看到的接口。
xv6 内核提供了 Unix 内核传统上所提供的服务和系统调用的一个子集。
图~\ref{fig:api}
列出了 xv6 的全部系统调用。

本章其余部分概述了 xv6 的服务——进程、内存、文件描述符、管道和文件系统——并通过代码片段及对\indextext{shell}（Unix 的命令行用户界面）如何使用这些服务的讨论加以说明。Shell 对系统调用的使用方式体现了这些接口设计的精心之处。

Shell 是一个普通程序，它从用户处读取命令并执行。
Shell 是用户程序而非内核的一部分，这一事实体现了系统调用接口的威力：shell 并无任何特殊之处。
这也意味着 shell 很容易被替换；因此，
现代 Unix 系统提供了多种
shell 供选择，各有其用户界面
和脚本功能。
xv6 的 shell 是 Unix Bourne shell 核心功能的简单实现。

xv6 shell 的实现可在
\fileref{user/sh.c} 中找到。（\textcolor{blue}{链接}是指向
\url{https://github.com/mit-pdos/xv6-riscv/} 上相关 xv6 源代码的超链接，具体编号参照
\href{xv6-src-booklet.pdf}{xv6-src-booklet.pdf} 中的页码和行号，类似于 Lions 的 UNIX 第六版注释~\cite{lions}。建议先自行阅读源代码（在你喜欢的开发环境中、在 GitHub 上或在 PDF 阅读器中），再回来参阅本书。学完本书后，你应该能够不借助本书理解 xv6 源代码的每一行。


%% 
%% 	Processes and memory
%% 
\section{进程与内存}

xv6 进程由用户空间内存（指令、数据和栈）以及内核私有的每进程状态构成。
xv6 对进程进行
\indextext{分时}处理：它在等待执行的进程集合之间透明地切换可用的 CPU。
当进程未在执行时，xv6 保存该进程的 CPU 寄存器，
并在下次运行该进程时恢复这些寄存器。
内核为每个进程关联一个进程标识符，即
\indexcode{PID}。

\begin{figure}[t]
\center
\begin{tabular}{ll}
{\bf System call} & {\bf Description} \\
\midrule
int fork() & 创建一个进程，返回子进程的 PID。 \\
int exit(int status) & 终止当前进程；状态报告给 wait()。不返回。 \\
int wait(int *status) & 等待子进程退出；退出状态存入 *status；返回子进程 PID。 \\
int kill(int pid) & 终止进程 PID。返回 0，出错返回 -1。 \\
int getpid() & 返回当前进程的 PID。 \\
int pause(int n) & 暂停 n 个时钟节拍。 \\
int exec(char *file, char *argv[]) & 加载文件并以参数执行；仅在出错时返回。 \\
char *sbrk(int n) & 将进程内存增加 n 个零字节。返回新内存的起始地址。 \\
int open(char *file, int flags) & 打开文件；flags 指示读/写方式；返回文件描述符 fd。 \\
int write(int fd, char *buf, int n) & 从 buf 向文件描述符 fd 写入 n 字节；返回 n。 \\
int read(int fd, char *buf, int n) & 将 n 字节读入 buf；返回已读字节数；文件末尾返回 0。 \\
int close(int fd) & 释放打开的文件描述符 fd。 \\
int dup(int fd) & 返回一个指向与 fd 相同文件的新文件描述符。\\
int pipe(int p[]) & 创建管道，将读/写文件描述符存入 p[0] 和 p[1]。 \\
int chdir(char *dir) & 改变当前目录。 \\
int mkdir(char *dir) & 创建新目录。 \\
int mknod(char *file, int, int) & 创建设备文件。 \\
int fstat(int fd, struct stat *st) & 将打开文件的信息存入 *st。 \\
int link(char *file1, char *file2) & 为文件 file1 创建另一个名称（file2）。 \\
int unlink(char *file) & 删除一个文件。 \\
\end{tabular}
\caption{Xv6 系统调用。除非另有说明，这些调用在无错误时返回 0，出错时返回 -1。}
\label{fig:api}
\end{figure}

进程可以使用
\indexcode{fork}
系统调用创建一个新进程。
\lstinline{fork}
为新进程提供调用进程内存的精确副本：\lstinline{fork} 将调用进程的指令、数据和
栈复制到新进程的内存中。
\lstinline{fork}
在原进程和新进程中都会返回。
在原进程中，\lstinline{fork} 返回新进程的
PID。
在新进程中，\lstinline{fork} 返回零。
原进程和新进程通常分别称为
\indextext{父}进程和
\indextext{子}进程。

例如，考虑以下用 C 语言~\cite{kernighan} 编写的程序片段：
\begin{lstlisting}[]
int pid = fork();
if(pid > 0){
  printf("parent: child=%d\n", pid);
  pid = wait((int *) 0);
  printf("child %d is done\n", pid);
} else if(pid == 0){
  printf("child: exiting\n");
  exit(0);
} else {
  printf("fork error\n");
}
\end{lstlisting}
\indexcode{exit}
系统调用使调用进程停止执行并释放内存和打开文件等资源。
\lstinline{exit} 接受一个整数状态参数，
按惯例 0 表示成功，1 表示失败。
\indexcode{wait}
系统调用返回当前进程的一个已退出（或被杀死）子进程的 PID，并将子进程的退出状态复制到传给 \lstinline{wait} 的地址；如果调用者的任何子进程
都尚未退出，
\indexcode{wait}
会等待某个子进程退出。
如果调用者没有子进程，\lstinline{wait} 立即
返回 -1。
如果父进程不关心子进程的退出状态，可以向
\lstinline{wait} 传入地址 0。

在上面的示例中，输出行
\begin{lstlisting}[]
parent: child=1234
child: exiting
\end{lstlisting}
的顺序可能有所不同（甚至交错出现），这取决于父进程还是子进程先执行到
\indexcode{printf}
调用。子进程退出后，父进程的
\indexcode{wait}
返回，导致父进程打印
\begin{lstlisting}[]
parent: child 1234 is done
\end{lstlisting}
虽然子进程以父进程内存的副本为起点，但父子进程使用各自独立的内存和寄存器：
在一个进程中修改变量不会影响另一个进程。例如，当
\lstinline{wait}
的返回值存入
\lstinline{pid}
在父进程中，
子进程中的变量
\lstinline{pid}
不会改变。子进程中
\lstinline{pid}
的值仍为零。

\indexcode{exec}
系统调用
用从文件系统中存储的文件加载的新内存映像替换调用进程的内存。
该文件必须具有特定格式，用于指定文件的哪部分存放指令、哪部分是数据、从哪条指令
开始执行等。xv6
使用 ELF 格式，第~\ref{CH:MEM}~章将对此进行更详细的讨论。
通常该文件是对程序源代码进行编译的结果。
当
\indexcode{exec}
成功执行时，它不会返回到调用程序；
相反，从文件中加载的指令会从 ELF 头中声明的入口点开始执行。
\lstinline{exec}
接受两个参数：包含可执行文件的文件名和一个字符串参数数组。
例如：
\begin{lstlisting}[]
char *argv[3];

argv[0] = "echo";
argv[1] = "hello";
argv[2] = 0;
exec("/bin/echo", argv);
printf("exec error\n");
\end{lstlisting}
该片段用以参数列表
\lstinline{echo}
\lstinline{hello} 运行的
\lstinline{/bin/echo}
程序实例替换调用程序。
大多数程序会忽略参数数组的第一个元素，按惯例
它是程序的名称。

xv6 shell 使用上述调用代表用户运行程序。Shell 的主体结构很简单，参见
\lstinline{main}
\lineref{user/sh.c:/main/}。
主循环通过
\indexcode{getcmd} 从用户读取一行输入，
然后调用
\lstinline{fork}，
创建 shell 进程的副本。父进程调用
\lstinline{wait}，
而子进程运行命令。例如，若用户
向 shell 输入了
``\lstinline{echo hello}''，
\lstinline{runcmd}
将以
``\lstinline{echo hello}''
为参数被调用。
\lstinline{runcmd}
\lineref{user/sh.c:/runcmd/}
执行实际命令。对于
``\lstinline{echo hello}''，
它将调用
\lstinline{exec}
\lineref{user/sh.c:/exec.ecmd/}。
若
\lstinline{exec}
成功，子进程将执行来自
\lstinline{echo}
的指令而非
\lstinline{runcmd}。
\lstinline{echo}
在某个时刻会调用
\lstinline{exit}，
这将导致父进程从
\lstinline{wait}
中的
\lstinline{main}
\lineref{user/sh.c:/main/} 返回。

你可能会疑惑为何
\indexcode{fork}
和
\indexcode{exec}
没有合并为一个调用；后面我们将看到，
shell 在实现 I/O 重定向时正是利用了二者的分离。
为了避免先创建进程副本然后立即用 \lstinline{exec} 替换它的浪费，
操作系统内核通过使用写时复制等虚拟内存技术（参见第~\ref{CH:PGFAULTS}~节）优化了
\lstinline{fork}
在此场景下的实现。

xv6 隐式分配大部分用户空间内存：
\indexcode{fork}
分配子进程复制父进程内存所需的内存，
\indexcode{exec}
分配足够容纳可执行文件的内存。
需要在运行时申请更多内存的进程（例如用于
\indexcode{malloc}）
可以调用
\lstinline{sbrk(n)}
将其数据内存增加
\lstinline{n} 个零
字节；
\indexcode{sbrk}
返回新内存的起始地址。

%% 
%% 	I/O and File descriptors
%% 
\section{I/O 与文件描述符}

\indextext{文件描述符}
是一个小整数，代表一个由内核管理的对象，
进程可以对其进行读写操作。
进程可以通过打开文件、目录或设备，或者创建管道，或者复制已有
描述符来获取文件描述符。
为简便起见，我们通常将文件描述符
所指向的对象称为"文件"；
文件描述符接口屏蔽了
文件、管道和设备之间的差异，使它们都呈现为字节流的形式。
我们将输入和输出统称为 \indextext{I/O}。

在内部，xv6 内核将文件描述符用作每进程表的索引，
因此每个进程拥有从零开始的私有文件描述符空间。
按惯例，进程从文件描述符 0（标准输入）读取数据，
向文件描述符 1（标准输出）写输出，
向文件描述符 2（标准错误）写错误消息。
如我们将看到的，shell 利用这一约定来实现 I/O 重定向
和管道。Shell 确保自身始终有三个文件描述符
处于打开状态
\lineref{user/sh.c:/open..console/}，
这三个描述符默认指向控制台。

\lstinline{read}
和
\lstinline{write}
系统调用分别从文件描述符所指向的打开文件中读取字节和向其写入字节。
调用
\lstinline{read(fd},
\lstinline{buf},
\lstinline{n)}
从文件描述符
\lstinline{fd}
中最多读取
\lstinline{n}
字节，将其复制到
\lstinline{buf}
中，并返回实际读取的字节数。
每个指向文件的文件描述符
都关联有一个偏移量。
\lstinline{read}
从当前文件偏移处读取数据，然后将
偏移量前进已读取的字节数：
后续的
\lstinline{read}
调用将返回第一次
\lstinline{read} 所返回字节之后的数据。
当没有更多字节可读时，
\lstinline{read}
返回零以表示文件末尾。

调用
\lstinline{write(fd},
\lstinline{buf},
\lstinline{n)}
将
\lstinline{buf}
中的
\lstinline{n}
字节写入文件描述符
\lstinline{fd}
并返回已写入的字节数。
只有在发生错误时才会写入少于
\lstinline{n}
字节。
与
\lstinline{read}
类似，
\lstinline{write}
从当前文件偏移处写入数据，然后将
偏移量前进已写入的字节数：
每次
\lstinline{write}
从上次结束的位置继续写入。

以下程序片段（构成
\lstinline{cat}
程序的核心）将数据从标准输入
复制到标准输出。如果发生错误，则向标准错误写入消息。
\begin{lstlisting}[]
char buf[512];
int n;

for(;;){
  n = read(0, buf, sizeof buf);
  if(n == 0)
    break;
  if(n < 0){
    fprintf(2, "read error\n");
    exit(1);
  }
  if(write(1, buf, n) != n){
    fprintf(2, "write error\n");
    exit(1);
  }
}
\end{lstlisting}
该代码片段中需要注意的重要一点是，
\lstinline{cat}
不知道它是从文件、控制台还是管道读取数据，
同样也不知道它是向控制台、文件还是其他设备输出数据。
文件描述符的使用以及文件描述符 0
为输入、文件描述符 1 为输出的约定，使得
\lstinline{cat} 的实现非常简单。

\lstinline{close}
系统调用
释放一个文件描述符，使其可供未来的
\lstinline{open}、
\lstinline{pipe}
或
\lstinline{dup}
系统调用（见下文）重新使用。
新分配的文件描述符
始终是当前进程中编号最小的未使用
描述符。

文件描述符与
\indexcode{fork}
的交互使得 I/O 重定向很容易实现。
\lstinline{fork}
在复制父进程内存的同时也复制其文件描述符表，
因此子进程与父进程拥有完全相同的打开文件。
系统调用
\indexcode{exec}
替换调用进程的内存，但保留其文件表。
这一行为使得 shell 可以通过 fork、在子进程中关闭并重新打开选定的文件描述符，
再调用 \lstinline{exec} 运行新程序来实现 \indextext{I/O 重定向}。
以下是 shell 执行命令
\lstinline{cat}
\lstinline{<}
\lstinline{input.txt} 时运行代码的简化版本：
\begin{lstlisting}[]
char *argv[2];

argv[0] = "cat";
argv[1] = 0;
if(fork() == 0) {
  close(0);
  open("input.txt", O_RDONLY);
  exec("cat", argv);
}
\end{lstlisting}
子进程关闭文件描述符 0 后，
\lstinline{open}
保证将该文件描述符分配给新打开的
\lstinline{input.txt}：
0 将是最小的可用文件描述符。
\lstinline{cat}
随后以文件描述符 0（标准输入）指向
\lstinline{input.txt} 的状态执行。
该操作序列不会改变父进程的文件描述符，
因为它只修改了子进程的描述符。

xv6 shell 中 I/O 重定向的代码正是以这种方式工作的
\lineref{user/sh.c:/case.REDIR/}。
请注意，在代码的这一点上，shell 已经 fork 了子 shell，
\lstinline{runcmd}
将调用
\lstinline{exec}
来加载新程序。

\lstinline{open} 的第二个参数是一组以位表示的标志，用于控制 \lstinline{open}
的行为。可能的取值定义在文件控制（fcntl）头文件
\linerefs{kernel/fcntl.h:/RDONLY/,/TRUNC/} 中：
\lstinline{O_RDONLY}、
\lstinline{O_WRONLY}、
\lstinline{O_RDWR}、
\lstinline{O_CREATE} 和
\lstinline{O_TRUNC}，
分别指示 \lstinline{open} 以只读、只写、读写方式打开文件，
在文件不存在时创建文件，
以及将文件截断为零长度。

现在应该清楚为什么
\lstinline{fork}
和
\lstinline{exec}
作为独立调用是有益的了：在两次调用之间，shell 有机会
重定向子进程的 I/O，而不干扰主 shell 的 I/O 设置。
可以设想一个假想的组合调用
\lstinline{forkexec}，
但用这样的调用来实现 I/O 重定向
似乎很笨拙。
Shell 可以在调用 \lstinline{forkexec} 前修改自身的 I/O
设置（随后再撤销这些修改）；或者
\lstinline{forkexec} 将 I/O
重定向指令作为参数；
或者（最不理想的方案）让每个类似 \lstinline{cat} 的程序
自己处理 I/O 重定向。

虽然
\lstinline{fork}
复制了文件描述符表，但底层的文件偏移量在父子进程之间是共享的。
考虑以下示例：
\begin{lstlisting}[]
if(fork() == 0) {
  write(1, "hello ", 6);
  exit(0);
} else {
  wait(0);
  write(1, "world\n", 6);
}
\end{lstlisting}
在该片段结束时，附加到文件描述符 1
的文件将包含数据
\lstinline{hello}
\lstinline{world}。
父进程中的
\lstinline{write}
（由于
\lstinline{wait}
的存在，它只在子进程完成后才运行）
从子进程的
\lstinline{write}
结束处继续写入。
这一行为有助于从 shell 命令序列（如
\lstinline{(echo}
\lstinline{hello};
\lstinline{echo}
\lstinline{world)}
\lstinline{>output.txt}）产生顺序输出。

\lstinline{dup}
系统调用复制一个现有的文件描述符，
返回一个指向相同底层 I/O 对象的新描述符。
两个文件描述符共享同一个偏移量，就像
\lstinline{fork}
复制的文件描述符那样。
这是向文件写入
\lstinline{hello}
\lstinline{world}
的另一种方式：
\begin{lstlisting}[]
fd = dup(1);
write(1, "hello ", 6);
write(fd, "world\n", 6);
\end{lstlisting}

如果两个文件描述符是通过一系列
\lstinline{fork}
和
\lstinline{dup}
调用从同一个原始文件描述符派生而来的，则它们共享同一个偏移量。
否则，即使两个文件描述符是由对同一文件的
\lstinline{open}
调用产生的，它们也不共享偏移量。
\lstinline{dup}
使 shell 能够实现如下命令：
\lstinline{ls}
\lstinline{existing-file}
\lstinline{non-existing-file}
\lstinline{>}
\lstinline{tmp1}
\lstinline{2>&1}。
其中
\lstinline{2>&1}
告诉 shell 将命令的文件描述符 2 设为描述符 1 的副本。
已存在文件的名称和不存在文件的错误消息都将出现在文件
\lstinline{tmp1}
中。
xv6 shell 不支持对错误文件描述符进行 I/O 重定向，但现在你已经知道如何实现它了。

文件描述符是一种强大的抽象，
因为它隐藏了所连接对象的细节：
向文件描述符 1 写入的进程，可能是在向文件、控制台等设备或管道写入数据。
%% 
%% 	Pipes
%% 
\section{管道}

\indextext{管道}
是一个小型内核缓冲区，以一对
文件描述符的形式暴露给进程，一个用于读取，一个用于写入。
向管道一端写入数据后，该数据即可从管道另一端读取。
管道为进程间通信提供了一种途径。

以下示例代码将
\lstinline{wc}
程序的标准输入连接到管道的读端来运行该程序。
\begin{lstlisting}[]
int p[2];
char *argv[2];

argv[0] = "wc";
argv[1] = 0;

pipe(p);
if(fork() == 0) {
  close(0);
  dup(p[0]);
  close(p[0]);
  close(p[1]);
  exec("/bin/wc", argv);
} else {
  close(p[0]);
  write(p[1], "hello world\n", 12);
  close(p[1]);
}
\end{lstlisting}
程序调用
\lstinline{pipe}，
创建一个新管道并将读写文件描述符记录在数组
\lstinline{p}
中。
\lstinline{fork}
之后，父子进程都拥有指向管道的文件描述符。
子进程调用 \lstinline{close} 和 \lstinline{dup}，使文件描述符
零指向管道的读端，
然后关闭
\lstinline{p}
中的文件描述符，
再调用 \lstinline{exec} 运行
\lstinline{wc}。
当
\lstinline{wc}
从其标准输入读取时，实际上是从管道读取。
父进程关闭管道的读端，
向管道写入数据，
然后关闭写端。

如果没有数据可读，对管道的
\lstinline{read}
调用将等待数据被写入，或者等到所有指向写端的
文件描述符都被关闭；
在后一种情况下，
\lstinline{read}
将返回 0，就像到达了数据文件的末尾一样。
\lstinline{read}
会阻塞直到不可能有新数据到来，
这正是上面的例子中子进程在执行
\lstinline{wc}
之前必须关闭管道写端的重要原因：如果
\lstinline{wc}
的某个文件描述符指向管道的写端，
\lstinline{wc}
将永远看不到文件末尾。

xv6 shell 以类似上述代码
\lineref{user/sh.c:/case.PIPE/}
的方式实现管道，如
\lstinline{grep fork sh.c | wc -l}。
子进程创建一个管道将管道线的左端与右端连接起来，然后为管道线的左端调用
\lstinline{fork}
和
\lstinline{runcmd}，
为右端调用
\lstinline{fork}
和
\lstinline{runcmd}，并等待两者完成。
管道线的右端可能是一个本身包含管道的命令（例如，
\lstinline{a}
\lstinline{|}
\lstinline{b}
\lstinline{|}
\lstinline{c)}，
它本身会 fork 两个新子进程（一个用于
\lstinline{b}，
一个用于
\lstinline{c}）。
因此，shell 可能会
创建一棵进程树。该树的叶节点是命令，
内部节点是等待左右子进程完成的进程。

% In principle, one could have the interior nodes run the left end of a
% pipeline, but doing so correctly would complicate the
% implementation. Consider making just the following modification:
% change \lstinline{sh.c} to not fork for \lstinline{p->left} and run
% \lstinline{runcmd(p->left)} in the interior process. Then, for
% example, \lstinline{echo hi | wc} won't produce output, because when
% \lstinline{echo hi} exits in \lstinline{runcmd}, the interior process
% exits and never calls fork to run the right end of the pipe.  This
% incorrect behavior could be fixed by not calling \lstinline{exit} in
% \lstinline{runcmd} for interior processes, but this fix complicates
% the code: now \lstinline{runcmd} needs to know if it's in an interior
% process or not.  Complications also arise when not forking for
% \lstinline{runcmd(p->right)}.  For example, with just that
% modification, \lstinline{sleep 10 | echo hi} will immediately print
% ``hi' and a new prompt, instead of after 10 seconds; this happens because \lstinline{echo} runs
% immediately and exits, not waiting for \lstinline{sleep} to finish.
% Since the goal of the \lstinline{sh.c} is to be as simple as possible,
% it doesn't try to avoid creating interior processes.

管道看起来可能并不比临时文件更强大：
管道线
\begin{lstlisting}[]
echo hello world | wc
\end{lstlisting}
可以不用管道实现为
\begin{lstlisting}[]
echo hello world >/tmp/xyz; wc </tmp/xyz
\end{lstlisting}
然而，在这种情况下，管道相比临时文件至少有三个优势。
第一，管道会自动清理自身；
而使用文件重定向，shell 必须在完成后小心地删除
\lstinline{/tmp/xyz}。
第二，管道可以传递任意长度的数据流，而文件重定向需要磁盘上有足够的
空闲空间来存储所有数据。
第三，管道允许管道线各阶段并行执行，
而文件方式要求第一个程序完成后第二个才能开始。
% Fourth, if you are implementing inter-process communication,
% pipes' blocking reads and writes are more efficient
% than the non-blocking semantics of files.
%% 
%% 	File system
%% 
\section{文件系统}

xv6 文件系统提供数据文件（包含未经解释的字节数组）
和目录（包含对数据文件和其他目录的命名引用）。
这些目录构成一棵树，从称为
\indextext{根}的特殊目录开始。
像
\lstinline{/a/b/c}
这样的\indextext{路径}
指的是根目录
\lstinline{/}
中目录
\lstinline{a}
下目录
\lstinline{b}
中名为
\lstinline{c}
的文件或目录。
不以
\lstinline{/}
开头的路径相对于调用进程的
\indextext{当前目录}进行求值，
当前目录可以通过
\lstinline{chdir}
系统调用更改。
以下两段代码都打开同一个文件
（假设所涉及的目录都存在）：
\begin{lstlisting}[]
chdir("/a");
chdir("b");
open("c", O_RDONLY);

open("/a/b/c", O_RDONLY);
\end{lstlisting}
第一段代码将进程的当前目录更改为
\lstinline{/a/b}；
第二段代码既不引用也不更改进程的当前目录。

有几个系统调用可以创建新文件和目录：
\lstinline{mkdir}
创建新目录，
\lstinline{open}
带
\lstinline{O_CREATE}
标志创建新数据文件，
\lstinline{mknod}
创建新设备文件。
以下示例展示了这三种情况：
\begin{lstlisting}[]
mkdir("/dir");
fd = open("/dir/file", O_CREATE|O_WRONLY);
close(fd);
mknod("/console", 1, 1);
\end{lstlisting}
\lstinline{mknod}
创建一个指向设备的特殊文件。
与设备文件关联的是主设备号和次设备号
（\lstinline{mknod} 的两个参数），
它们唯一标识一个内核设备。
当进程之后打开设备文件时，内核
将
\lstinline{read}
和
\lstinline{write}
系统调用转向内核设备实现，
而不是传递给文件系统。

文件名与文件本身是不同的；
同一个底层文件（称为
\indextext{inode}）
可以有多个名称，
称为
\indextext{链接}。
每个链接由目录中的一个条目组成；
该条目包含一个文件名和对
inode 的引用。
inode 保存了文件的
\indextext{元数据}，
包括其类型（文件、目录或设备）、
大小、
文件内容在磁盘上的位置，
以及指向该文件的链接数。

\lstinline{fstat}
系统调用
从文件描述符所指向的 inode 中检索信息。
它填充一个
\lstinline{struct}
\lstinline{stat}，
该结构体在
\lstinline{stat.h} \fileref{kernel/stat.h}
中定义如下：
\begin{lstlisting}[]
#define T_DIR     1   // Directory
#define T_FILE    2   // File
#define T_DEVICE  3   // Device

struct stat {
  int dev;     // File system's disk device
  uint ino;    // Inode number
  short type;  // Type of file
  short nlink; // Number of links to file
  uint64 size; // Size of file in bytes
};
\end{lstlisting}

\lstinline{link}
系统调用创建另一个文件系统名称，
指向与现有文件相同的 inode。
以下片段创建了一个同时命名为
\lstinline{a}
和
\lstinline{b} 的新文件。
\begin{lstlisting}[]
open("a", O_CREATE|O_WRONLY);
link("a", "b");
\end{lstlisting}
对
\lstinline{a}
的读写与对
\lstinline{b}
的读写效果相同。
每个 inode 由唯一的
\textit{inode}
\textit{编号}标识。
在上述代码序列之后，可以通过检查
\lstinline{fstat}
的结果来确定
\lstinline{a}
和
\lstinline{b}
指向相同的底层内容：两者都将返回相同的 inode 编号
（\lstinline{ino}），
且
\lstinline{nlink}
计数将被设置为 2。

\lstinline{unlink}
系统调用从文件系统中删除一个名称。
只有当文件的链接计数为零且没有文件描述符引用它时，文件的 inode 和保存其内容的磁盘空间才会被释放。
因此，在上一段代码序列中添加
\begin{lstlisting}[]
unlink("a");
\end{lstlisting}
后，inode 和文件内容仍可通过
\lstinline{b}
访问。
此外，
\begin{lstlisting}[]
fd = open("/tmp/xyz", O_CREATE|O_RDWR);
unlink("/tmp/xyz");
\end{lstlisting}
是创建一个无名临时 inode 的惯用方式，该 inode 将在进程关闭
\lstinline{fd}
或退出时被清理。

Unix 将文件工具作为用户级程序提供，可从 shell 调用，
例如
\lstinline{mkdir}、
\lstinline{ln}
和
\lstinline{rm}。
这种设计允许任何人通过添加新的用户级程序来扩展命令行界面。事后看来这个方案显而易见，
但与 Unix 同时代设计的其他系统往往将此类命令内置于
shell（并将 shell 内置于内核）。

一个例外是
\lstinline{cd}，
它被内置于 shell
\lineref{user/sh.c:/if.cmd.0..==..c./} 中。
\lstinline{cd}
必须更改 shell 本身的当前工作目录。如果
\lstinline{cd}
作为普通命令运行，则 shell 会 fork 一个子
进程，子进程运行
\lstinline{cd}，
\lstinline{cd}
会更改
\textit{子}进程的
工作目录。父进程（即
shell）的工作目录不会改变。
%% 
%% 	Real world
%% 
\section{现实世界}

Unix 将"标准"文件
描述符、管道以及操作它们的便捷 shell 语法有机结合，是编写
通用可复用程序的重大进步。
这一理念催生了"软件工具"文化，这也是
Unix 强大功能和广泛普及的重要原因，
而 shell 是第一种所谓的"脚本语言"。
Unix 系统调用接口至今仍在
BSD、Linux 和 macOS 等系统中延续。

Unix 系统调用接口已通过可移植操作系统接口（POSIX）标准进行了标准化。
xv6
\textit{不}
符合 POSIX 标准：它缺少许多系统调用（包括
\lstinline{lseek} 等基本调用），
且其提供的许多系统调用也与标准有所不同。
xv6 的主要目标是
简洁明了，同时提供简单的类 Unix 系统调用接口。
一些人通过添加更多系统调用和简单的
C 库扩展了 xv6，以便运行基本的 Unix 程序。然而，现代内核
提供的系统调用和内核服务远多于
xv6，例如支持网络、窗口系统、用户级线程、
多种设备驱动等。现代内核持续快速演进，提供了许多超出 POSIX 标准的功能。

Unix 通过一套文件名和文件描述符接口统一了对多种资源（文件、
目录和设备）的访问。
这一理念可以扩展到更多类型的资源；
一个典型示例是 Plan 9~\cite{Presotto91plan9}，
它将"资源即文件"的
概念应用于
网络、图形等领域。
然而，大多数源自 Unix 的操作系统并未
沿循这条路线。

文件系统和文件描述符一直是强大的抽象机制。
尽管如此，操作系统接口还有其他模型。
Unix 的前身 Multics
以一种使文件存储看起来像内存的方式进行了抽象，
产生了截然不同风格的接口。
Multics 设计的复杂性对 Unix 的设计者产生了直接影响，他们的目标是构建更简单的系统。
% XXX can we cut this, since its point is the same as the next paragraph?
% An operating system interface that went out of fashion
% decades ago but has recently returned is the idea of a virtual machine monitor.
% Such systems provide a superficially different interface from xv6,
% but the basic concepts are still the same:
% a virtual machine, like a process, consists of some memory and
% one or more register sets;
% the virtual machine has access to one large file called
% a virtual disk instead of a file system;
% virtual machines send messages to each other
% and the outside world using virtual network devices
% instead of pipes or files.

Xv6 does not provide a notion of users or of protecting
one user from another; in Unix terms, all xv6 processes
run as root.

This book examines how xv6 implements its Unix-like interface,
but the ideas and concepts apply to more than just Unix.
Any operating system must multiplex processes onto
the underlying hardware, isolate processes from each
other, and provide mechanisms for controlled
inter-process communication.
After studying xv6, you should be able to
look at other, more complex operating systems
and see the concepts underlying xv6 in those systems as well.

%% 
\section{Exercises}
%%

\begin{enumerate}

\item Write a program that uses UNIX system calls to
``ping-pong'' a byte between two processes over a pair
of pipes, one for each direction. Measure the
program's performance, in exchanges per second.

\end{enumerate}
